\chapter{Hertz's Mechanics Without Force}
\label{hm}

By the 1890s, several physicists --- Kirchhoff and Mach among them
--- had grown uneasy with the status of \emph{force} as a
fundamental concept of mechanics. Force is defined operationally
through $F=ma$, but that same equation is routinely used to
\emph{introduce} force, which makes the law read uncomfortably like
a definition dressed up as an empirical statement. In his
posthumously published \emph{Die Prinzipien der Mechanik in neuem
Zusammenhange dargestellt} (1894), Heinrich Hertz proposed to remove
the discomfort by removing force altogether: he set out to build the
whole of mechanics from just three primitive notions --- space,
time, and mass --- with no independent concept of force at
all~\cite{hertz}.

\section{The principle of the straightest path}

Consider $N$ mass points with positions $\mathbf x_i\in\mathbb R^3$
and masses $m_i$, and equip the $3N$-dimensional configuration space
$Q$ with the natural mass-weighted (kinetic-energy) metric,
\begin{equation}
  ds^2 \;=\; \sum_{i=1}^N m_i\, d\mathbf x_i\cdot d\mathbf x_i,
  \qquad
  T \;=\; \frac12\left(\frac{ds}{dt}\right)^{\!2}.
  \label{hertz-metric}
\end{equation}
Hertz's masses are never free in the naive sense: they are always
tied together by \emph{rigid, workless connections} (idealized
mechanical linkages that transmit force but do no work), which
restrict the accessible configurations to some submanifold
$\mathcal C\subset Q$. His fundamental postulate --- the
\emph{principle of the straightest path}
--- is then simply:

\begin{quote}
\emph{A system free of "force" moves along the straightest possible
path compatible with its constraints, at constant speed.}
\end{quote}

"Straightest possible" is to be read literally, in the geometric
sense: the representative point traces a \emph{geodesic} of the
metric~\eqref{hertz-metric} restricted to $\mathcal C$, i.e.\ a
curve whose acceleration (in the induced metric) is purely normal to
$\mathcal C$ and is entirely accounted for by the constraint
reaction. Formally, if the constraints are written $g_a(\mathbf
x)=0$, $a=1,\dots,k$, the equations of motion are
\begin{equation}
  m_i\,\ddot{\mathbf x}_i \;=\; \sum_a \lambda_a\,
  \frac{\partial g_a}{\partial \mathbf x_i},
  \label{hertz-eom}
\end{equation}
with the multipliers $\lambda_a$ fixed by demanding $g_a(\mathbf
x(t))\equiv 0$. There is nothing on the right-hand side of
\eqref{hertz-eom} except constraint reactions: no applied-force
term appears anywhere, by construction. Constant speed
($|ds/dt|=\mathrm{const}$, i.e.\ constant kinetic energy) follows
automatically, since ideal constraints do no work.

\section{Hidden masses, or where "force" goes}

Of course, ordinary bodies \emph{do} appear to exert forces on one
another --- gravity, contact forces, and so on all look like genuine
interactions between visible masses that are not rigidly connected
to anything. Hertz's resolution was to posit \emph{hidden masses}:
invisible mass points, rigidly linked to the visible ones (and, in
his more speculative moments, identified with a mechanical
aether), whose only role is to curve the constraint submanifold
$\mathcal C$ in just the right way. If the full configuration space
splits as $Q = Q_{\mathrm{vis}}\times Q_{\mathrm{hid}}$ and the
\emph{complete} system --- visible plus hidden --- moves along a
geodesic of $\mathcal C$, then projecting that geodesic motion down
onto $Q_{\mathrm{vis}}$ alone generically produces a curve that is
\emph{not} geodesic in the visible sector by itself: extra terms
appear in its equation of motion that are, from the point of view of
an observer with access only to $Q_{\mathrm{vis}}$, indistinguishable
from the effect of an applied force~\cite{hertz}.

This is a genuinely striking claim: on Hertz's account, force is not
absent from experience, only absent from the fundamental ontology.
Every observed force law would, in principle, be explained by some
specific (if permanently unobservable) arrangement of hidden mass
and rigid linkage. The same basic move --- explain an apparent force
as the shadow, on a submanifold, of pure geometric (geodesic) motion
in a larger space --- reappears twice in twentieth-century physics
in far more successful forms: in Kaluza--Klein theory, where the
electromagnetic force on a charged particle is recovered as geodesic
motion in a five-dimensional space, and, far more importantly, in
general relativity, where gravity itself --- which Newton had to
treat as a force acting at a distance --- becomes geodesic motion of
a free particle through curved four-dimensional spacetime, with no
"force" required at all. Hertz did not anticipate curved spacetime;
his geometry is Euclidean and the curvature is manufactured entirely
by hidden extra coordinates. But the structural resemblance to
Einstein's later, physically successful geometrization of gravity is
close enough that Hertz's book is still read today chiefly for that
reason, rather than as a working theory in its own
right.

\section{A second, more conservative geometrization}

It is worth noting that ordinary, force-\emph{having} mechanics can
also be written as a geodesic principle, without any hidden masses
at all, via Jacobi's form of the principle of least action. For a
conservative system with potential $V(\mathbf x)$ and fixed total
energy $E$, the true trajectories are exactly the geodesics of the
\emph{Jacobi metric}
\begin{equation}
  ds_J^2 \;=\; 2\bigl(E-V(\mathbf x)\bigr)\, ds^2 ,
  \label{jacobi-metric}
\end{equation}
a conformal rescaling of the mass metric~\eqref{hertz-metric} by
the local kinetic energy. So there are, in the end, two quite
different ways to eliminate force from the description of a system
governed by a potential: enlarge the configuration space with hidden
coordinates and keep the metric flat (Hertz), or keep the visible
configuration space as it is and let the metric itself become curved
and energy-dependent (Jacobi). Both reproduce exactly the same
observable trajectories. Neither is more "true" than the other; they
are, in Hertz's own vocabulary, different \emph{images} of the same
phenomena. That a concept as apparently basic as force can be made
to evaporate --- twice over, in two structurally different ways ---
without changing a single observable prediction is the real content
of this section: it shows "force" to be a feature of a chosen
representation of mechanics, not a fact about the world that every
correct representation must contain.

Hertz's own programme did not survive contact with its critics.
Almost every contemporary reviewer, Boltzmann included, objected
that positing a fresh, permanently unobservable mechanism of hidden
masses for every distinct force law (gravitation, electromagnetism,
\ldots) bought elegance at the price of testability, and the
research programme was essentially abandoned within a generation. It
survives today less as working physics than as an early, and
unusually explicit, demonstration that mechanics has more than one
internally consistent axiomatic foundation --- a lesson later put to
far more productive use.
